% ---------------------------------------------------------------- Telnet
\section{Telnet}
\label{sec:telnet}

The remote shell in this project started as a Telnet server and client
(commit \texttt{4b08310}). Telnet showed how a remote terminal works, and also
why it cannot be used on a real network: everything it sends, passwords
included, can be read by anyone on the path. The rsh programs keep Telnet's
idea of a terminal session over TCP and add the missing security. This section
describes Telnet and the first implementation; Section~\ref{sec:rshdesign}
describes how it became rsh.

\subsection{The protocol}

Telnet is defined in RFC~854 (1983). A client opens a TCP connection to port
23 and both sides then exchange a stream of bytes in the format of the
\emph{Network Virtual Terminal} (NVT): 7-bit ASCII text, with every line ending
in \texttt{CR LF}. There are no packets or headers of Telnet's own; the
protocol lives entirely in the byte stream.

Control information is mixed into the data with an escape byte,
\texttt{IAC} (``interpret as command'', 255). The byte after \texttt{IAC} is a
command. A data byte that really is 255 is sent as \texttt{IAC IAC}.

\begin{table}[H]
\centering
\small
\begin{tabularx}{\textwidth}{@{}r p{3.4cm} X@{}}
\toprule
Code & Name & Meaning \\
\midrule
240 & \texttt{SE}   & end of a subnegotiation \\
241 & \texttt{NOP}  & no operation \\
242--249 & \texttt{DM BRK IP AO}\newline \texttt{AYT EC EL GA} & data mark, break, interrupt process, abort output, are you there, erase character, erase line, go ahead \\
250 & \texttt{SB}   & start of a subnegotiation: \texttt{IAC SB} \textit{option} \ldots\ \texttt{IAC SE} \\
251 & \texttt{WILL} & sender wants to enable an option on its side \\
252 & \texttt{WONT} & sender refuses to enable, or disables, an option on its side \\
253 & \texttt{DO}   & sender asks the other side to enable an option \\
254 & \texttt{DONT} & sender asks the other side not to use, or to stop, an option \\
255 & \texttt{IAC}  & interpret the next byte as a command; \texttt{IAC IAC} is a data byte 255 \\
\bottomrule
\end{tabularx}
\caption{Telnet commands (RFC 854).}
\end{table}

\textbf{Option negotiation} (RFC~855) lets the two ends agree on extra
features. Each option is a number, and a request takes three bytes, for
example \texttt{IAC DO 1} (``please echo''). The answer is either the
agreeing form (\texttt{WILL 1}) or the refusing form (\texttt{WONT 1}), and a
side must not answer a request for a state it is already in, which prevents
endless loops. Either side may always refuse, so a minimal implementation that
refuses everything still works with any client: both ends then fall back to
the plain NVT.

\begin{table}[H]
\centering
\small
\begin{tabular}{@{}r l l l@{}}
\toprule
Number & Option & RFC & Purpose \\
\midrule
0  & Binary transmission & 856  & send 8-bit data instead of NVT ASCII \\
1  & Echo                & 857  & the side that \texttt{WILL ECHO} echoes the characters it receives \\
3  & Suppress go-ahead   & 858  & full-duplex use, no \texttt{GA} after each turn \\
24 & Terminal type       & 1091 & the client names its terminal, e.g.\ \texttt{xterm} \\
31 & Window size (NAWS)  & 1073 & the client reports columns and rows \\
34 & Linemode            & 1184 & edit lines locally and send whole lines \\
\bottomrule
\end{tabular}
\caption{Common Telnet options.}
\end{table}

Telnet has no login protocol of its own. The server simply runs a
\texttt{login} program, which prints \texttt{login:} and \texttt{Password:} and
reads the answers from the same byte stream as everything else. Nothing is
encrypted, the server proves nothing about who it is, and the client proves
who it is only by typing a password across the network. For these reasons
Telnet was replaced by SSH, first released in 1995 and standardised in
RFC~4251--4254 (2006).

\subsection{Our Telnet implementation}

The first version (\texttt{telnet/}, about 260 lines of Rust, standard library only) had a
multi-threaded server on TCP port 2323 and an interactive client. Instead of a
real shell the server ran a small command interpreter with \texttt{help},
\texttt{echo <text>}, \texttt{time} and \texttt{quit}.

Both sides passed every received chunk through the same \texttt{Parser}, a
state machine that separates data bytes from \texttt{IAC} commands
(Figure~\ref{fig:telnetfsm}). It refuses every option, answering \texttt{DO}
with \texttt{WONT} and \texttt{WILL} with \texttt{DONT}, and skips
subnegotiations. Outgoing data goes through \texttt{escape}, which doubles each
255 byte.

\begin{figure}[H]
\centering
\begin{tikzpicture}[
  >={Stealth[length=2mm]}, node distance=2.1cm and 2.9cm, font=\small,
  st/.style={draw, rounded corners, minimum width=1.9cm, minimum height=0.75cm, fill=gray!8},
  lab/.style={font=\footnotesize, fill=white, inner sep=1.5pt},
]
  \node[st] (data) {\texttt{Data}};
  \node[st, right=of data] (iac) {\texttt{Iac}};
  \node[st, right=of iac] (opt) {\texttt{Option(cmd)}};
  \node[st, below=1.9cm of iac] (sub) {\texttt{Sub}};
  \node[st, right=of sub] (subiac) {\texttt{SubIac}};

  \draw[->] (data) to[loop left] node[lab, left]{other: emit byte} (data);
  \draw[->] (data) to[bend left=18] node[lab, above]{255} (iac);
  \draw[->] (iac) to[bend left=18] node[lab, below]{255: emit 255; other: ignore} (data);
  \draw[->] (iac) -- node[lab, above]{\texttt{DO}, \texttt{DONT}, \texttt{WILL}, \texttt{WONT}} (opt);
  \draw[->] (opt) to[out=110, in=60, looseness=0.55] node[lab, above]{option byte: reply \texttt{WONT}/\texttt{DONT}} (data);
  \draw[->] (iac) -- node[lab, right]{\texttt{SB}} (sub);
  \draw[->] (sub) to[loop below] node[lab, below]{other} (sub);
  \draw[->] (sub) to[bend left=15] node[lab, above]{255} (subiac);
  \draw[->] (subiac) to[bend left=15] node[lab, below]{other} (sub);
  \draw[->] (subiac) to[out=200, in=-80, looseness=1.1] node[lab, below left]{\texttt{SE}} (data);
\end{tikzpicture}
\caption{The Telnet input parser of the first version (\texttt{protocol.rs}).}
\label{fig:telnetfsm}
\end{figure}

\Needspace{16\baselineskip}
\begin{lstlisting}[caption={Option refusal and escaping in the Telnet parser (\texttt{protocol.rs}).}]
State::Iac => match b {
    IAC => { data.push(IAC); State::Data }           // IAC IAC is a data byte 255
    DO | DONT | WILL | WONT => State::Option(b),
    SB => State::Sub,
    _ => State::Data,                                  // NOP, GA, ...: ignored
},
State::Option(cmd) => {
    match cmd {
        DO => replies.extend([IAC, WONT, b]),          // "I will not"
        WILL => replies.extend([IAC, DONT, b]),        // "please do not"
        _ => {}                                        // WONT, DONT need no answer
    }
    State::Data
}
\end{lstlisting}

\subsection{A Telnet session on the wire}
\label{sec:telnetwire}

To show what Telnet exposes, the first version was rebuilt from the repository
history and a session was recorded through a TCP relay that copies every byte
(\texttt{report/tools/relay.py}). The client first sent three option requests
of the kind a Telnet client sends when it connects, then typed four lines. This is the complete
traffic, in order:

% generated by tools/gen_tables.py from data/ -- do not edit
\begin{lstlisting}[style=plain]
server -> client   63 B  "Welcome to the Rust telnet server\r\nType 'help' for commands\r\n> "
client -> server    9 B  ff fd 03 ff fb 18 ff fb 1f
                         = IAC DO SUPPRESS-GO-AHEAD, IAC WILL TERMINAL-TYPE, IAC WILL NAWS
server -> client    9 B  ff fc 03 ff fe 18 ff fe 1f
                         = IAC WONT SUPPRESS-GO-AHEAD, IAC DONT TERMINAL-TYPE, IAC DONT NAWS
client -> server   15 B  "login student\r\n"
server -> client   26 B  "Unknown command: login\r\n> "
client -> server   22 B  "echo my-password-123\r\n"
server -> client   19 B  "my-password-123\r\n> "
client -> server    6 B  "time\r\n"
server -> client   40 B  "Seconds since Unix epoch: 1794883707\r\n> "
client -> server    6 B  "quit\r\n"
server -> client    9 B  "Goodbye\r\n"
\end{lstlisting}


The server refused all three options, exactly as in
Figure~\ref{fig:telnetfsm}. Everything else is readable as it is. The line
that carried the ``password'' looks like this on the wire:

\begin{table}[H]
\centering
% generated by tools/gen_tables.py from data/ -- do not edit
{\setlength{\tabcolsep}{1.6pt}\footnotesize
\begin{tabular}{@{}lcccccccccccccccccccccc@{}}
\toprule
wire byte & \texttt{65} & \texttt{63} & \texttt{68} & \texttt{6f} & \texttt{20} & \texttt{6d} & \texttt{79} & \texttt{2d} & \texttt{70} & \texttt{61} & \texttt{73} & \texttt{73} & \texttt{77} & \texttt{6f} & \texttt{72} & \texttt{64} & \texttt{2d} & \texttt{31} & \texttt{32} & \texttt{33} & \texttt{0d} & \texttt{0a} \\
as text & \texttt{e} & \texttt{c} & \texttt{h} & \texttt{o} & \texttt{\textvisiblespace} & \texttt{m} & \texttt{y} & \texttt{-} & \texttt{p} & \texttt{a} & \texttt{s} & \texttt{s} & \texttt{w} & \texttt{o} & \texttt{r} & \texttt{d} & \texttt{-} & \texttt{1} & \texttt{2} & \texttt{3} & \texttt{\textbackslash r} & \texttt{\textbackslash n} \\
\bottomrule
\end{tabular}}

\caption{The bytes of \texttt{echo my-password-123} as captured: each byte is the ASCII code of the character typed.}
\label{tab:telnetpw}
\end{table}

Anyone who can see the packets (another host on the same Wi-Fi, a router, the
ISP) reads the password directly. They can also change the bytes or inject new
commands, and the client has no way to tell whether it is talking to the real
server at all.

\subsection{What an encrypted Telnet needs}

Turning Telnet into a safe remote shell needs four things that Telnet does not
have. Each became part of rsh:

\begin{enumerate}
\item \textbf{Server authentication}: the client must know it reached the right
  server. rsh gives each server an RSA host key and keeps a
  \texttt{known\_hosts} list.
\item \textbf{A shared secret}: both ends need a key that the observer does not
  have. The client makes a random session key and sends it encrypted with the
  host's public key.
\item \textbf{Confidentiality}: every byte after the key exchange is encrypted
  with a stream cipher keyed by the session key.
\item \textbf{Client authentication without a password on the wire}: the client
  signs the session key with its private key; the server checks the signature
  against \texttt{authorized\_keys}.
\end{enumerate}

\noindent A fifth requirement, \emph{integrity} (detecting changed bytes), is
not met by this version; Section~\ref{sec:cryptolimits} shows what that allows.
